% part: intuitionistic-logic
% chapter: propositions-as-types
% section: sequent-natural-deduction

\documentclass[../../../include/open-logic-section]{subfiles}

\begin{document}

\olfileid{int}{pty}{snd}

\olsection{સિક્વન્ટ-શૈલીનું સ્વાભાવિક નિગમન}

સ્વાભાવિક નિગમનની એવી !!{derivation} હોય
કે તેની !!{undischarged} ધારણાઓ~$\Gamma$
અને નિષ્કર્ષ~$!A$ હોય, તો
$\Gamma \Sequent !A$ લખીએ;
$\Gamma$ ખાલી હોય, તો $\Sequent !A$ લખીએ.

$\Gamma \cup \{!A_1, \dots, !A_n\}$ માટે
$\Gamma, !A_1, \dots, !A_n$ અને
$\Gamma \cup \Delta$ માટે $\Gamma, \Delta$ લખીએ.

નોંધો કે $\Gamma \Sequent !A \land !B$ હોય,
એટલે !!{undischarged} ધારણાઓ~$\Gamma$ અને
અંતિમ !!{formula}~$!A \land !B$ ધરાવતી
!!a{derivation} હોય, તો સૌથી નીચે
$\Elim{\land}$ લાગુ કરીને એ જ
!!{undischarged} ધારણાઓ અને નિષ્કર્ષ~$!A$
ધરાવતી !!a{derivation} મળે છે.
એટલે $\Gamma \Sequent !A \land !B$
હોય, તો $\Gamma \Sequent !A$ છે.
\sourcecorrection{OLMD-295}{મૂળ આ ખુલાસાની શરૂઆતમાં બંને સંયોજ્યો $!A$ લખે છે, પરંતુ સારવાક્ય અને નિયમવૃક્ષમાં $!A$ તથા $!B$ છે. ખુલાસાના બંને સ્થાનોને એ જ સામાન્ય સંયોજન સાથે સુસંગત કર્યા છે.}
\begin{prooftree}
  \Axiom$\Gamma \fCenter !A \land !B$
  \RightLabel{$\Elim{\land}$}
  \UnaryInf$\Gamma \fCenter !A$
  \DisplayProof\qquad\bottomAlignProof
  \Axiom$\Gamma \fCenter !A \land !B$
  \RightLabel{$\Elim{\land}$}
  \UnaryInf$\Gamma \fCenter !B$
\end{prooftree}
$\Elim{\land}$ ચિહ્ન સ્વાભાવિક નિગમનના
એ જ નામના નિયમ સાથેનો સંબંધ સૂચવે છે.

એ જ રીતે $\Gamma, !A \Sequent !B$ હોય,
એટલે !!{undischarged} ધારણાઓ $\Gamma, !A$
અને અંતિમ !!{formula}~$!B$ ધરાવતી
!!a{derivation} હોય એમ ધારો.
$\Intro{\lif}$ નિયમ લગાવીએ, તો
!!{undischarged} ધારણાઓ~$\Gamma$
અને અંતિમ !!{formula}~$!A \lif !B$
ધરાવતી !!a{derivation} મળે છે,
એટલે $\Gamma \Sequent !A \lif !B$.
આ રીતે ધારણાઓની !!{discharge}
વધુ સ્પષ્ટ થાય છે તે નોંધો.
\begin{prooftree}
  \Axiom $\Gamma, !A \fCenter !B$
  \RightLabel{$\Intro{\lif}$}
  \UnaryInf $\Gamma \fCenter !A \lif !B$
\end{prooftree}

બીજા નિયમોમાંથી પણ આ રીતે નિષ્કર્ષો
મેળવી શકીએ. તે આમ સ્પષ્ટ થાય છે:
\begin{gather*}
  \Axiom $\Gamma \fCenter !A$
  \Axiom $\Delta \fCenter !B$
  \RightLabel{$\Intro{\land}$}
  \BinaryInf $\Gamma,\Delta \fCenter !A \land !B$
  \DisplayProof\\
  \Axiom $\Gamma \fCenter !A \land !B$
  \RightLabel{$\Elim{\land}_1$}
  \UnaryInf $\Gamma \fCenter !A$
  \DisplayProof
  \qquad
  \Axiom $\Gamma \fCenter !A \land !B$
  \RightLabel{$\Elim{\land}_2$}
  \UnaryInf $\Gamma \fCenter !B$
  \DisplayProof
  \\
  \Axiom $\Gamma \fCenter !A$
  \RightLabel{$\Intro{\lor}_1$}
  \UnaryInf $\Gamma \fCenter !A \lor !B$
  \DisplayProof\qquad
  \Axiom $\Gamma \fCenter !B$
  \RightLabel{$\Intro{\lor}_2$}
  \UnaryInf $\Gamma \fCenter !A \lor !B$
  \DisplayProof\\
  \Axiom $\Gamma \fCenter !A \lor !B$
  \Axiom $\Delta, !A \fCenter !C$
  \Axiom $\Delta', !B \fCenter !C$
  \RightLabel{$\Elim{\lor}$}
  \TrinaryInf $\Gamma, \Delta, \Delta' \fCenter !C$
  \DisplayProof
  \\
  \Axiom $\Gamma, !A \fCenter !B$
  \RightLabel{$\Intro{\lif}$}
  \UnaryInf $\Gamma \fCenter !A \lif !B$
  \DisplayProof
  \qquad
  \Axiom $\Delta \fCenter !A \lif !B$
  \Axiom $\Gamma \fCenter !A$
  \RightLabel{$\Elim{\lif}$}
  \BinaryInf $\Gamma, \Delta \fCenter !B$
  \DisplayProof
  \\
  \Axiom $\Gamma \fCenter \lfalse$
  \RightLabel{$\FalseInt$}
  \UnaryInf $\Gamma \fCenter !A$
  \DisplayProof
\end{gather*}

કોઈ પણ ધારણા પોતે~$!A$માંથી $!A$ની
!!a{derivation} છે; એટલે હંમેશાં
$!A \Sequent !A$ છે.

\begin{prooftree}
  \AxiomC{$ $}
  \UnaryInfC{$!A \fCenter !A$}
\end{prooftree}

આ બધા નિયમોને સ્વાભાવિક નિગમનની કઈ
!!{derivation}s અસ્તિત્વ ધરાવે છે તે
વિશેનું કલન ગણી શકાય. તેને સ્વાભાવિક
નિગમનની નોંધની વૈકલ્પિક રીત પણ
ગણી શકાય, જેમાં દરેક પગલે ફક્ત
!!{derive}d થયેલું !!{formula} જ
નહીં, પણ જે !!{undischarged} ધારણાઓમાંથી
તે !!{derive}d થયું હોય તે પણ નોંધાય છે.

\begin{prooftree}
  \Axiom$!A \fCenter !A$
  \UnaryInf $!A \fCenter !A \lor (!A \lif \lfalse)$
  \Axiom $!B \fCenter !B$
  \BinaryInf$!A, !B \fCenter \lfalse$
  \UnaryInf$!B \fCenter !A \lif \lfalse$
  \UnaryInf$!B \fCenter !A \lor (!A \lif \lfalse)$
  \Axiom$!B \fCenter !B$
  \BinaryInf$!B \fCenter \lfalse$
  \UnaryInf$\fCenter !B \lif \lfalse$
\end{prooftree}
અહીં $!B$ એ
$(!A \lor (!A \lif \lfalse))\lif \lfalse$નું
ટૂંકું રૂપ છે.
\sourcecorrection{OLMD-296}{મૂળના અંતિમ વૃક્ષમાં પૂર્વાંગના સૂત્ર $!B$ પછી એક વધારાનો અનુસરણ-સંયોજક તથા ચાર વધારાના ખુલ્લા કૌંસ છે. તેમને કાઢ્યાં છે; તેથી દરેક પગલાનો સંદર્ભ આપેલી $!B$ ધારણા છે અને અંતિમ નિષ્કર્ષ એ ધારણાથી અસત્ય મળે છે તે જ કહે છે.}

\end{document}
